\subsection{Strategies for Encouraging Collaboration Among AI Systems}
\label{sec:5.5.2}
Three levers do most of the work, and each has a failure mode worth stating alongside it.

Open tooling lets researchers outside a lab inspect and replicate what the lab claims, in a field where claims about a model's behavior are otherwise taken on faith \autocite{wolf2020transformers}. Openness accumulated as a side effect of release practice does less than openness pointed deliberately at alignment questions that benchmark scores do not capture \autocite{alignmentforum2018}. Multi-agent research platforms are the specific case: they let norm formation and adversarial drift be studied under conditions closer to deployment than a single-agent benchmark allows \autocite{leibo2021meltingpot}.

Competitions pull toward whatever they measure. A leaderboard on a fixed benchmark says nothing about ethical behavior in either direction; a competition judged on a real problem the entrants chose themselves gives them a reason to weigh who benefits.

\begin{esbox}
\boxtitle{The IBM Watson AI XPRIZE, 2016-2021}
A five-year, \$5 million competition run by the XPRIZE Foundation with IBM, open to any team willing to apply AI to a "grand challenge" of its own choosing, with no fixed benchmark set for it. Finalists worked on malaria elimination, mental health treatment, and detecting human trafficking; Zzapp Malaria won the grand prize in 2021 \autocite{xprize2021watson}.

What is worth taking from it is the structure rather than the winner: judging teams on a problem they chose themselves is one way a competition can pull toward collaboration.

\end{esbox}
Standing consortia work on a longer clock and are the weakest of the three. Section~\ref{sec:8.4.2} gives the Partnership on AI its full treatment, including what a voluntary body can and cannot do, and section~\ref{sec:7.4} gives the argument that a consortium of this kind is itself a channel worth testing. The international version is section~\ref{sec:8.7.1}'s.

None of this works without people who are not AI researchers. Neuroscience, psychology, philosophy and political science each catch failure modes the others are structurally blind to, which is chapter~\ref{sec:8}'s opening argument rather than an afterthought here.
